\documentclass[11pt]{article}

\usepackage[T1]{fontenc}
\usepackage{amsmath,amssymb}
\usepackage{hyperref}

\title{H-Weighted Robust Polychorics}
\author{}
\date{\today}

\newcommand{\Rmat}{\widehat R^{(h)}}
\newcommand{\thetahat}{\widehat\theta^{(h)}}
\newcommand{\fhat}{\widehat f}

\begin{document}
\maketitle

\noindent
Reference: Welz, Mair, and Alfons, ``Robust Estimation of Polychoric
Correlation,'' \emph{Psychometrika} 91, 247--278, 2026,
doi: \href{https://doi.org/10.1017/psy.2025.10066}{10.1017/psy.2025.10066}.
The local paper copy is
\path{resources/papers/robust_ordinal/welz_etal_2025_robust_polychoric_correlation.pdf}.

\section*{Bivariate model}

Let $X_a$ and $X_b$ be ordinal variables with $K_a$ and $K_b$ categories. The
polychoric model introduces latent variables $(Z_a,Z_b)$ with standard
bivariate normal distribution and correlation $\rho$. Thresholds
\[
  -\infty = \tau_{a0} < \tau_{a1} < \cdots < \tau_{a,K_a-1}
  < \tau_{a,K_a} = \infty
\]
and similarly for variable $b$ map latent values to ordinal cells. For a cell
$j=(r,s)$,
\[
  p_j(\theta) =
  \Pr\{\tau_{a,r-1} < Z_a \le \tau_{a,r},
       \tau_{b,s-1} < Z_b \le \tau_{b,s}; \rho\},
  \qquad
  \theta = (\tau_a,\tau_b,\rho).
\]
With observed cell count $n_j$ and pair sample size $n$, write
$\fhat_j=n_j/n$ and
\[
  t_j(\theta) = \frac{\fhat_j}{p_j(\theta)}, \qquad
  s_j(\theta) = \nabla_\theta \log p_j(\theta).
\]

\section*{H-score definition}

The h-weighted robust polychoric estimator solves
\[
  \Psi_h(\theta)
  =
  \sum_j p_j(\theta) h\{t_j(\theta)\} s_j(\theta)
  = 0,
\]
or equivalently minimizes the associated minimum-disparity objective when the
chosen $h$ has a corresponding loss $\Phi$:
\[
  Q_h(\theta)
  =
  \sum_j p_j(\theta)
  \Phi\!\left(\frac{\fhat_j}{p_j(\theta)}\right).
\]
Ordinary ML is the special case $h(t)=t$. The Welz--Mair--Alfons hard-cap
choice is
\[
  h(t) = \min(t,k),
\]
with $k=c+1$ in their notation. The implementation also keeps smooth capped
variants for optimizer experiments, but the hard cap is the paper-facing
robust polychoric definition.

\section*{Pair-local versus SEM moments}

For a single bivariate table, the pair-local robust polychoric is the
correlation component of
\[
  \thetahat_{ab}
  =
  \arg\min_{\tau_a,\tau_b,\rho}
  Q_h(\tau_a,\tau_b,\rho),
  \qquad
  \widehat\rho_{ab,h} =
  \rho(\thetahat_{ab}).
\]
This is the direct Welz--Mair--Alfons object: thresholds are nuisance
parameters local to the pair.

For ordinal SEM sample statistics, magmaan uses one shared threshold vector per
observed variable and one robust correlation per pair. The shared-threshold
composite target is
\[
  \widehat\eta_h =
  \arg\min_{\{\tau_a\},\{\rho_{ab}\}}
  \sum_{a<b}\sum_{j\in ab}
  p_{ab,j}(\tau_a,\tau_b,\rho_{ab})
  \Phi\!\left(
    \frac{\widehat f_{ab,j}}
         {p_{ab,j}(\tau_a,\tau_b,\rho_{ab})}
  \right).
\]
The robust polychoric matrix is then
\[
  \Rmat_{aa}=1,
  \qquad
  \Rmat_{ab}=\Rmat_{ba}=\widehat\rho_{ab,h}.
\]
Because pairwise robust correlations are not automatically positive definite,
the implementation treats positive-definiteness repair as explicit diagnostics:
none, error, ridge, or shrinkage toward the identity.

\section*{Sandwich object}

Let $\psi_i(\eta)$ be the stacked h-score contribution for row $i$ in the
shared-threshold composite estimating equation. At the solution,
\[
  A_h =
  \left.
  \frac{\partial}{\partial \eta^\top}
  \left\{
    \frac{1}{n}\sum_i \psi_i(\eta)
  \right\}
  \right|_{\widehat\eta_h},
  \qquad
  B_h =
  \frac{1}{n}\sum_i
  \psi_i(\widehat\eta_h)\psi_i(\widehat\eta_h)^\top .
\]
The moment influence rows are
\[
  IF_i = -\psi_i(\widehat\eta_h) A_h^{-\top},
\]
and the ordinal moment covariance used downstream is
\[
  \Gamma_h = \frac{1}{n}\sum_i IF_i IF_i^\top .
\]
The moment order is magmaan's ordinal order: all thresholds first, followed by
lower-triangle polychorics by columns.

\section*{Current magmaan access}

There is no separate friendly R helper named only for ``robust polychoric
matrix.'' The matrix is already obtainable from the methods-developer data
builder without fitting a SEM:
\begin{verbatim}
d <- magmaan_core$data_ordinal_stats_from_raw(
  list(X),
  robust = "h_weighted",
  h_kind = "wma_hard_cap",
  h_k = 1.5
)
Rh <- d$R[[1]]
\end{verbatim}
For data frames, use:
\begin{verbatim}
d <- magmaan_core$data_ordinal_stats_from_df(
  df, model, robust = "h_weighted"
)
Rh <- d$R[[1]]
\end{verbatim}

The corresponding C++ hooks are:
\begin{flushleft}\footnotesize
\path{data::pairwise_ordinal_stats_h_weighted_from_integer_data()}\\
shared-threshold SEM-grade matrix.

\path{data::fit_ordinal_pair_joint_h_weighted()}\\
\path{data::fit_ordinal_pair_rho_h_weighted()}\\
direct bivariate work.
\end{flushleft}

\end{document}
